% ANCHOR: after_label def:c5-Oeps
\begin{remark}[Physics notation: sources and renormalised insertions]
\begin{itemize}
\item $\omega=\omega_{1234}(x)\,dx^1\wedge dx^2\wedge dx^3\wedge dx^4$ is a compactly supported source and $O=\int d^4x\,\omega_{1234}(x)\,(\text{density})$ couples it to a composite operator; $\delta$ keeps the first order in the source, and the Weyl weight $w_O$ is the mass dimension of the composite operator.
\item $O(\varepsilon)=O-\partial_\delta I^{CT}[\varepsilon]$ is the bare composite operator with its one-loop operator counterterms: operator renormalisation and mixing, including mixing with the identity $\int\omega$.
\item $O[\infty]$ is the renormalised insertion: the $\varepsilon\to0$ limit of the connected graphs with one vertex $O(\varepsilon)$, vertices from $I-I^{CT}[\varepsilon]$, lines $\hbar P(\varepsilon,\infty)$ and the background fields on the legs.
\item $(d_{CE}+Q)\,O[\infty]$ is the action on the insertion of the free BRST differential and of the constant superconformal ghosts.
\end{itemize}
\end{remark}

% ANCHOR: end_of_section 5.1
\begin{remark}[Physics notation: operator counterterms and BRST invariance]
Proposition \ref{prop:c5-thmB}: the graphs with the insertion off the loop are renormalised by the action counterterm $I^{CT}[\varepsilon]$, and the operator counterterm is the $\varepsilon^{-1}$ and $\log\varepsilon$ part of the graphs with the insertion on the loop, whose superficial degree of divergence is $w_O$ minus the mass dimensions of the legs. The cocycle condition $Q\,O[\infty]=0$ is BRST invariance of the renormalised insertion. An output of the type $\int\omega\,(\partial_mc,A_m)$ is removed by a finite multiple of the gluon mass operator $\int\omega\,(A_m,A_m)$, since
\[
Q\int\omega\,(A_m,A_m)=-2\int\omega\,(\partial_mc,A_m),
\]
from the arrow $c\mapsto-\partial_mc$ of $Q$.
\end{remark}

% ANCHOR: end_of_section 5.2
\begin{remark}[Physics notation: colour contractions, vertices, logarithms]
For $\mathfrak{su}(N)$, $\hbar(-,-)^{-1}=\tfrac12\sum_a(-igT^a)\otimes(-igT^a)$, and with $X=-igX^aT^a$ Lemma \ref{lem:c5-colour} reads
\[
[-,-]\ \text{on}\ (-,-)^{-1}:\ f^{abc}\delta^{ab}=0,\qquad (-,-)\ \text{on}\ (-,-)^{-1}=N^2-1,\qquad ([X_1,-],[X_2,-])\ \text{on}\ (-,-)^{-1}=-\operatorname{tr}_{\mathfrak g}(\mathrm{ad}_{X_1}\mathrm{ad}_{X_2})=g^2N\,X_1^aX_2^a .
\]
The cubic vertex of Lemma \ref{lem:c5-vertices} is the gluon--scalar--scalar vertex of $-\tfrac1{16}(D_mA^{ij},D_mA_{ij})=-\tfrac14\sum_I(D_mA_I,D_mA_I)$, and the quartic one the quartic part of the bosonic action. Lemma \ref{lem:c5-log}: the $\log\varepsilon$ coefficient of a logarithmically divergent one-loop graph is $-\frac1{16\pi^2}$ times the angular average of the top-degree part of its numerator; for the basic bubble
\[
\int\frac{d^4p}{(2\pi)^4}\,\frac{e^{-\varepsilon|p|^2}}{|p|^2}\,\frac{e^{-\varepsilon|p+p_1|^2}}{|p+p_1|^2}\ \supset\ -\frac{\log\varepsilon}{16\pi^2},\qquad \Big(\frac{1-e^{-|x-y|^2/4\varepsilon}}{4\pi^2|x-y|^2}\Big)^2\ \supset\ -\frac{\log\varepsilon}{16\pi^2}\,\delta(x-y).
\]
\end{remark}

% ANCHOR: end_of_section 5.3
\begin{remark}[Physics notation: single letters]
In components, for $X\in\mathfrak g$, $w\in W_{\mathbb R}^\perp$, $v\in V$ and $\zeta=(\zeta^\alpha{}_i,\bar\zeta^{\dot\alpha i})\in\Sigma^*\otimes\mathfrak g$:
\[
\int\omega\,(X,\iota_wA)=\int\omega\,w_I(X,A_I)=\tfrac14\int\omega\,w^{ij}(X,A_{ij}),\qquad \int\omega\,(X,\iota_v(A|_V))=\int\omega\,v_m(X,A_m),
\]
\[
\int\omega\,(\zeta,\lambda)=\int\omega\,\big[(\lambda_\alpha{}^i,\zeta^\alpha{}_i)-(\bar\lambda_{\dot\alpha i},\bar\zeta^{\dot\alpha i})\big],\qquad \int\omega\,(X,c).
\]
For $\mathfrak{su}(N)$, $(X,A_I)=2\hbar X^aA^a_I$ and $(X,A_m)=2\hbar X^aA^a_m$. The letters are the elementary fields: the scalars $A^a_I$ in the $\mathbf 6$ of $SO(6)$ (as $A_{ij}$, the antisymmetric $\mathbf 6$ of $SU(4)$), the gluon $A^a_m$, the gauginos $\lambda_\alpha{}^i$ in $(\mathbf 2,\mathbf 1,\mathbf 4)$ and $\bar\lambda_{\dot\alpha i}$ in $(\mathbf 1,\mathbf 2,\bar{\mathbf 4})$ with sources $\zeta^\alpha{}_i$, $\bar\zeta^{\dot\alpha i}$, and the ghost $c$. Proposition \ref{prop:c5-letters}: a letter lies off every loop, the tadpoles carry $f^{abc}\delta^{bc}=0$, the other loops are renormalised by $I^{CT}[\varepsilon]$, and the operator counterterm vanishes, $O(\varepsilon)=O$.
\end{remark}

% ANCHOR: after_label def:c5-bilinear
\begin{remark}[Physics notation: $\mathbf{20'}$ and Konishi]
With $\phi_I=A^a_IT^a$, i.e. $A_I=-ig\,\phi_I$, one has $(A_I,A_J)=4\hbar\operatorname{tr}(\phi_I\phi_J)=2\hbar\,\phi^a_I\phi^a_J$ and
\[
\int\omega\,(\iota_{w_1}A,\iota_{w_2}A)=w_{1I}w_{2J}\int\omega\,(A_I,A_J)=\tfrac1{16}\,w_1^{ij}w_2^{kl}\int\omega\,(A_{ij},A_{kl})=4\hbar\,w_{1I}w_{2J}\int\omega\operatorname{tr}(\phi_I\phi_J).
\]
Under $\mathrm{Spin}(6)\cong SU(4)$, $\mathrm{Sym}^2\mathbf 6=\mathbf{20'}\oplus\mathbf 1$. The bracket of (\ref{eq:c5-decomp}) is the $\mathbf{20'}$ ($[0,2,0]$ of $SU(4)$); in $SO(6)$ and in $SU(4)$ form, with $\phi_{ij}=\phi_I\Sigma^I_{ij}$ and $\phi^{ij}=\tfrac12\varepsilon^{ijkl}\phi_{kl}$,
\[
4\hbar\,w_{1I}w_{2J}\int\omega\operatorname{tr}\big(\phi_I\phi_J-\tfrac16\delta_{IJ}\phi_K\phi_K\big)=\tfrac\hbar4\,w_1^{ij}w_2^{kl}\int\omega\Big[\operatorname{tr}(\phi_{ij}\phi_{kl})-\tfrac1{12}\varepsilon_{ijkl}\operatorname{tr}(\phi^{mn}\phi_{mn})\Big].
\]
The trace part is $\tfrac16(w_1,w_2)$, $(w_1,w_2)=w_{1I}w_{2I}$, times the Konishi operator
\[
\int\omega\,(A|_{W_{\mathbb R}^\perp},A|_{W_{\mathbb R}^\perp})=\tfrac14\int\omega\,(A^{ij},A_{ij})=4\hbar\int\omega\operatorname{tr}(\phi_I\phi_I)=\hbar\int\omega\operatorname{tr}(\phi^{ij}\phi_{ij}).
\]
\end{remark}

% ANCHOR: after_label prop:c5-census
\begin{remark}[Physics notation: the insertion graphs]
In physics terms: (a) the self-contraction $\langle\phi^a_I\phi^b_J\rangle$ of the two scalars of the operator at coincident points; (b) one external gluon through a gluon--scalar--scalar vertex, with colour factor $f^{abc}\delta^{bc}=0$; (c) the seagull bubble through $-\tfrac14\sum_I([A_m,A_I],[A_m,A_I])$ and (d) the triangle through two gluon--scalar--scalar vertices, both with external gluons; (e) the bubble through the scalar potential $-\tfrac18\sum_{I,J}([A_I,A_J],[A_I,A_J])$ and (f) the triangle with a gluon exchanged between the two scalar lines, both with external scalars. With $w_O=2$, power counting leaves at most two external bosons.
\end{remark}

% ANCHOR: after_label thm:c5-Oeps
\begin{remark}[Physics notation: renormalised $\mathbf{20'}$ and Konishi operators]
In components,
\[
\partial_\delta I^{CT}[\varepsilon]=\mathrm{unit}\,\log\varepsilon\Big[\tfrac14\int\omega\,(w_1^{ij}A_{ij},w_2^{kl}A_{kl})-\tfrac18\,w_1^{ij}w_{2ij}\int\omega\,(A^{kl},A_{kl})\Big]+\frac{\hbar\dim\mathfrak g\;w_1^{ij}w_{2ij}}{32\pi^2\,\varepsilon}\int\omega,
\]
$\mathrm{unit}=\frac{\hbar\operatorname{tr}_{\mathfrak g}(\mathrm{ad}_X\mathrm{ad}_Y)}{16\pi^2(X,Y)}$. For $\mathfrak{su}(N)$ ($\mathrm{unit}=-\frac{g^2N}{32\pi^2}$, $\dim\mathfrak g=N^2-1$) every $\mathbf{20'}$ operator $O$ is multiplicatively renormalised,
\[
O(\varepsilon)=\Big(1+\frac{g^2N}{8\pi^2}\log\varepsilon\Big)O,
\]
and the Konishi operator $\int\omega\,(A|_{W_{\mathbb R}^\perp},A|_{W_{\mathbb R}^\perp})=2\hbar\int\omega\,\phi^a_I\phi^a_I$ mixes with the identity alone,
\[
\int\omega\,\phi^a_I\phi^a_I\ \longmapsto\ \Big(1-\frac{g^2N}{4\pi^2}\log\varepsilon\Big)\int\omega\,\phi^a_I\phi^a_I-\frac{3(N^2-1)}{8\pi^2\,\varepsilon}\int\omega,
\]
the last term subtracting the coincident-point value $\langle\phi^a_I\phi^a_I\rangle=\frac{6(N^2-1)}{16\pi^2\varepsilon}$ of graph (a).
\end{remark}
